\documentclass[11pt,a4paper]{article}
\usepackage[utf8]{inputenc}
\usepackage[T1]{fontenc}
\usepackage[spanish,es-minimal,es-noshorthands]{babel}
\usepackage{lmodern}
\usepackage[sfdefault,lining]{FiraSans}
\usepackage[a4paper,top=2.3cm,bottom=2.2cm,left=1.8cm,right=1.8cm,headheight=1.3cm,headsep=0.4cm,footskip=1cm]{geometry}
\usepackage{amsmath,amssymb}
\usepackage{icomma}
\usepackage{graphicx}
\usepackage[table]{xcolor}
\usepackage{tikz}
\usetikzlibrary{calc}
\usepackage{fontawesome5}
\usepackage{ccicons}
\usepackage{fancyhdr}
\usepackage{lastpage}
\usepackage{tasks}
\usepackage{needspace}
\usepackage[most]{tcolorbox}
\usepackage{array}
\usepackage[hidelinks]{hyperref}
\AtBeginDocument{\hypersetup{pdftitle={Ficha de ejercicios de la unidad 2: Álgebra. Ecuaciones e inecuaciones. 1º Bach.},pdfsubject={Matemáticas I 1º Bach.},pdfkeywords={polinomios, Ruffini, teorema del resto, factorización, fracciones algebraicas, ecuaciones, inecuaciones}}}

\definecolor{navy}{HTML}{1B2748}
\definecolor{accent}{HTML}{E8603A}
\definecolor{teal}{HTML}{1F9E90}
\definecolor{soft}{HTML}{EEF1F8}
\definecolor{muted}{HTML}{5D6685}

\setlength{\parindent}{0pt}
\setlength{\parskip}{3pt}
\settasks{label=\alph*), label-format=\bfseries\color{navy}, label-width=1.3em, item-indent=1.8em, column-sep=0.6em, after-item-skip=5pt, before-skip=3pt, after-skip=4pt}

% ---------- capçalera i peu ----------
\newcommand{\llicencia}{\href{https://creativecommons.org/licenses/by-nc-sa/4.0/deed.es}{\ccbyncsa}\ \ Licencia Creative Commons Atribución-NoComercial-CompartirIgual 4.0 Internacional (CC BY-NC-SA 4.0)}
\pagestyle{fancy}
\fancyhf{}
\fancyhead[L]{\raisebox{-0.35\height}{\includegraphics[height=0.95cm]{logo-gva.png}}}
\fancyhead[C]{\small\color{muted}Ficha de ejercicios · Unidad 2: Álgebra. Ecuaciones e inecuaciones · 1\textsuperscript{o} Bach.}
\fancyhead[R]{\raisebox{-0.35\height}{\includegraphics[height=0.95cm]{logo-vermellar.png}}}
\fancyfoot[L]{\scriptsize\color{muted}\llicencia}
\fancyfoot[R]{\small\color{muted}\thepage}
\renewcommand{\headrulewidth}{0.4pt}
\renewcommand{\headrule}{\color{muted!40}\hrule height \headrulewidth}
\fancypagestyle{primera}{\fancyhf{}\renewcommand{\headrulewidth}{0pt}
  \fancyfoot[L]{\scriptsize\color{muted}\llicencia}
  \fancyfoot[R]{\small\color{muted}\thepage}}

% ---------- nivells i exercicis ----------
\newcommand{\estrelles}[1]{{\color{accent}\ifcase#1\or\faStar\faStar[regular]\faStar[regular]\or\faStar\faStar\faStar[regular]\or\faStar\faStar\faStar\fi}}
\newcounter{exer}
\newcommand{\ex}[1]{\par\needspace{4\baselineskip}\medskip\refstepcounter{exer}\noindent{\color{navy}\bfseries\theexer.}\ \ }

\newcommand{\seccio}[2]{\par\needspace{16\baselineskip}\bigskip
  \begin{tcolorbox}[colback=navy,colframe=navy,coltext=white,arc=2mm,boxrule=0pt,left=3mm,right=3mm,top=1.2mm,bottom=1.2mm]
  {\large\bfseries #1\quad #2}
  \end{tcolorbox}\vspace{-2pt}}

\newcommand{\exf}[1]{\par\medskip\refstepcounter{exer}\noindent{\color{navy}\bfseries\theexer.}\ \ }
\providecommand{\lt}{<}\providecommand{\gt}{>}
\newcolumntype{C}[1]{>{\centering\arraybackslash}p{#1}}
% quadre de recordatori (només a la fitxa d'exercicis)
\newcommand{\recorda}[1]{\par\smallskip\begin{tcolorbox}[colback=accent!7,colframe=accent!55,arc=2mm,boxrule=0.5pt,left=2.5mm,right=2.5mm,top=1mm,bottom=1mm]\small{\color{accent}\faLightbulb[regular]}\ \textbf{Recuerda:}\ #1\end{tcolorbox}}
\newcommand{\res}[1]{#1}\newcommand{\rest}[1]{#1}

\def\fmneed{11}% que el títol d'una secció no quede sol al final del full
% =====================================================================
% Capçalera FentMates de les fitxes: un color per a cada curs.
% Els colors són els de la pàgina inicial (LEVELS a assets/js/materials.js);
% el logo ve de _eines/logo/pdf/ (python3 _eines/logo/genera.py).
%
% Ús, al final del preàmbul de la unitat:
%   % =====================================================================
% Capçalera FentMates de les fitxes: un color per a cada curs.
% Els colors són els de la pàgina inicial (LEVELS a assets/js/materials.js);
% el logo ve de _eines/logo/pdf/ (python3 _eines/logo/genera.py).
%
% Ús, al final del preàmbul de la unitat:
%   % =====================================================================
% Capçalera FentMates de les fitxes: un color per a cada curs.
% Els colors són els de la pàgina inicial (LEVELS a assets/js/materials.js);
% el logo ve de _eines/logo/pdf/ (python3 _eines/logo/genera.py).
%
% Ús, al final del preàmbul de la unitat:
%   \input{../comu/fentmates.tex}
%   \fmcurs{1eso}
% i, al començament del document:
%   \capcalera{Matemàtiques}{1r ESO}{Fitxa d'exercicis · Unitat 2}{Divisibilitat}
% Opcional abans de \input: \def\fmneed{14} (espai mínim abans de cada secció).
% =====================================================================
\makeatletter
\def\fm@logodir{../../logo/pdf}
\providecommand{\fmneed}{10}
\def\fm@cursnom{}

% clau del curs, color i color del text damunt del color
\newcommand{\fm@curs}[3]{\@namedef{fm@color@#1}{#2}\@namedef{fm@text@#1}{#3}}
\fm@curs{1eso}{E6B84A}{navy}
\fm@curs{2eso}{3B82F6}{white}
\fm@curs{3eso}{FF7A50}{white}
\fm@curs{4eso-a}{A78BFA}{white}
\fm@curs{4eso-b}{6D28D9}{white}
\fm@curs{1batx-mi}{9B1C1C}{white}
\fm@curs{1batx-ccss}{EF6B6B}{white}
\fm@curs{2batx-mii}{1E6B47}{white}
\fm@curs{2batx-ccss}{2BB3A3}{white}

\newcommand{\fmcurs}[1]{%
  \@ifundefined{fm@color@#1}{\PackageError{fentmates}{Curs desconegut: #1}{}}{}%
  \edef\fm@tmp{\noexpand\definecolor{curs}{HTML}{\@nameuse{fm@color@#1}}%
               \noexpand\colorlet{curstext}{\@nameuse{fm@text@#1}}}\fm@tmp
  \colorlet{cursfosc}{curs!72!black}%
  \def\fm@logo{\fm@logodir/vertical-#1.pdf}%
  \def\fm@marca{\fm@logodir/marca-#1.pdf}}

% ---------- capçalera de la primera pàgina ----------
% #1 matèria (petit), #2 curs, #3 tipus de document i unitat, #4 títol
\newcommand{\capcalera}[4]{%
  \gdef\fm@cursnom{#2}%
  \par\noindent
  \begin{tcolorbox}[colback=curs,colframe=curs,boxrule=0pt,arc=2.2mm,width=3.5cm,height=2.15cm,
      valign=center,left=3mm,right=2mm,top=1mm,bottom=1mm,nobeforeafter,box align=center]
    \color{curstext}{\footnotesize #1}\par\vspace{1pt}{\fontsize{19}{21}\selectfont\bfseries #2\par}%
  \end{tcolorbox}\hspace{5mm}%
  \begin{minipage}[c]{\dimexpr\linewidth-3.5cm-5mm-2.9cm\relax}
    {\color{cursfosc}\small\bfseries\MakeUppercase{#3}}\par\vspace{2pt}
    {\color{navy}\fontsize{20}{23}\selectfont\bfseries #4\par}
  \end{minipage}\hfill
  \begin{minipage}[c]{2.6cm}\raggedleft\includegraphics[width=2.6cm]{\fm@logo}\end{minipage}\par
  \vspace{2pt}}

% ---------- barres de secció del color del curs ----------
\renewcommand{\seccio}[2]{\par\needspace{\fmneed\baselineskip}\bigskip
  \begin{tcolorbox}[colback=curs,colframe=curs,coltext=curstext,arc=2mm,boxrule=0pt,left=3mm,right=3mm,top=1.2mm,bottom=1.2mm]
  {\large\bfseries #1\quad #2}
  \end{tcolorbox}\vspace{-2pt}}

% ---------- estrelles de dificultat del color del curs (a Batxillerat no se'n posen) ----------
\@ifundefined{estrelles}{}{%
  \renewcommand{\estrelles}[1]{{\color{curs}\ifcase#1\or\faStar\faStar[regular]\faStar[regular]\or\faStar\faStar\faStar[regular]\or\faStar\faStar\faStar\fi}}}

% ---------- capçalera i peu de les pàgines ----------
% primera pàgina: logos institucionals al peu
\fancypagestyle{primera}{\fancyhf{}\renewcommand{\headrulewidth}{0pt}
  \fancyfoot[L]{\raisebox{-0.32\height}{\includegraphics[height=0.78cm]{logo-gva.png}}\hspace{5mm}%
    \raisebox{-0.32\height}{\includegraphics[height=0.78cm]{logo-vermellar.png}}}
  \fancyfoot[R]{\scriptsize\color{muted}\href{https://creativecommons.org/licenses/by-nc-sa/4.0/deed.ca}{\ccbyncsa}\ \ CC BY-NC-SA 4.0 · fentmates.org\qquad\small\thepage}}
% la resta: el símbol a l'esquerra i el curs a la dreta, amb el títol al centre com abans
\fancyhead[L]{\raisebox{-0.3\height}{\includegraphics[height=0.62cm]{\fm@marca}}}
\fancyhead[R]{\small\bfseries\color{cursfosc}\fm@cursnom}
\renewcommand{\headrule}{\color{curs!45}\hrule height \headrulewidth}
\makeatother

%   \fmcurs{1eso}
% i, al començament del document:
%   \capcalera{Matemàtiques}{1r ESO}{Fitxa d'exercicis · Unitat 2}{Divisibilitat}
% Opcional abans de \input: \def\fmneed{14} (espai mínim abans de cada secció).
% =====================================================================
\makeatletter
\def\fm@logodir{../../logo/pdf}
\providecommand{\fmneed}{10}
\def\fm@cursnom{}

% clau del curs, color i color del text damunt del color
\newcommand{\fm@curs}[3]{\@namedef{fm@color@#1}{#2}\@namedef{fm@text@#1}{#3}}
\fm@curs{1eso}{E6B84A}{navy}
\fm@curs{2eso}{3B82F6}{white}
\fm@curs{3eso}{FF7A50}{white}
\fm@curs{4eso-a}{A78BFA}{white}
\fm@curs{4eso-b}{6D28D9}{white}
\fm@curs{1batx-mi}{9B1C1C}{white}
\fm@curs{1batx-ccss}{EF6B6B}{white}
\fm@curs{2batx-mii}{1E6B47}{white}
\fm@curs{2batx-ccss}{2BB3A3}{white}

\newcommand{\fmcurs}[1]{%
  \@ifundefined{fm@color@#1}{\PackageError{fentmates}{Curs desconegut: #1}{}}{}%
  \edef\fm@tmp{\noexpand\definecolor{curs}{HTML}{\@nameuse{fm@color@#1}}%
               \noexpand\colorlet{curstext}{\@nameuse{fm@text@#1}}}\fm@tmp
  \colorlet{cursfosc}{curs!72!black}%
  \def\fm@logo{\fm@logodir/vertical-#1.pdf}%
  \def\fm@marca{\fm@logodir/marca-#1.pdf}}

% ---------- capçalera de la primera pàgina ----------
% #1 matèria (petit), #2 curs, #3 tipus de document i unitat, #4 títol
\newcommand{\capcalera}[4]{%
  \gdef\fm@cursnom{#2}%
  \par\noindent
  \begin{tcolorbox}[colback=curs,colframe=curs,boxrule=0pt,arc=2.2mm,width=3.5cm,height=2.15cm,
      valign=center,left=3mm,right=2mm,top=1mm,bottom=1mm,nobeforeafter,box align=center]
    \color{curstext}{\footnotesize #1}\par\vspace{1pt}{\fontsize{19}{21}\selectfont\bfseries #2\par}%
  \end{tcolorbox}\hspace{5mm}%
  \begin{minipage}[c]{\dimexpr\linewidth-3.5cm-5mm-2.9cm\relax}
    {\color{cursfosc}\small\bfseries\MakeUppercase{#3}}\par\vspace{2pt}
    {\color{navy}\fontsize{20}{23}\selectfont\bfseries #4\par}
  \end{minipage}\hfill
  \begin{minipage}[c]{2.6cm}\raggedleft\includegraphics[width=2.6cm]{\fm@logo}\end{minipage}\par
  \vspace{2pt}}

% ---------- barres de secció del color del curs ----------
\renewcommand{\seccio}[2]{\par\needspace{\fmneed\baselineskip}\bigskip
  \begin{tcolorbox}[colback=curs,colframe=curs,coltext=curstext,arc=2mm,boxrule=0pt,left=3mm,right=3mm,top=1.2mm,bottom=1.2mm]
  {\large\bfseries #1\quad #2}
  \end{tcolorbox}\vspace{-2pt}}

% ---------- estrelles de dificultat del color del curs (a Batxillerat no se'n posen) ----------
\@ifundefined{estrelles}{}{%
  \renewcommand{\estrelles}[1]{{\color{curs}\ifcase#1\or\faStar\faStar[regular]\faStar[regular]\or\faStar\faStar\faStar[regular]\or\faStar\faStar\faStar\fi}}}

% ---------- capçalera i peu de les pàgines ----------
% primera pàgina: logos institucionals al peu
\fancypagestyle{primera}{\fancyhf{}\renewcommand{\headrulewidth}{0pt}
  \fancyfoot[L]{\raisebox{-0.32\height}{\includegraphics[height=0.78cm]{logo-gva.png}}\hspace{5mm}%
    \raisebox{-0.32\height}{\includegraphics[height=0.78cm]{logo-vermellar.png}}}
  \fancyfoot[R]{\scriptsize\color{muted}\href{https://creativecommons.org/licenses/by-nc-sa/4.0/deed.ca}{\ccbyncsa}\ \ CC BY-NC-SA 4.0 · fentmates.org\qquad\small\thepage}}
% la resta: el símbol a l'esquerra i el curs a la dreta, amb el títol al centre com abans
\fancyhead[L]{\raisebox{-0.3\height}{\includegraphics[height=0.62cm]{\fm@marca}}}
\fancyhead[R]{\small\bfseries\color{cursfosc}\fm@cursnom}
\renewcommand{\headrule}{\color{curs!45}\hrule height \headrulewidth}
\makeatother

%   \fmcurs{1eso}
% i, al començament del document:
%   \capcalera{Matemàtiques}{1r ESO}{Fitxa d'exercicis · Unitat 2}{Divisibilitat}
% Opcional abans de \input: \def\fmneed{14} (espai mínim abans de cada secció).
% =====================================================================
\makeatletter
\def\fm@logodir{../../logo/pdf}
\providecommand{\fmneed}{10}
\def\fm@cursnom{}

% clau del curs, color i color del text damunt del color
\newcommand{\fm@curs}[3]{\@namedef{fm@color@#1}{#2}\@namedef{fm@text@#1}{#3}}
\fm@curs{1eso}{E6B84A}{navy}
\fm@curs{2eso}{3B82F6}{white}
\fm@curs{3eso}{FF7A50}{white}
\fm@curs{4eso-a}{A78BFA}{white}
\fm@curs{4eso-b}{6D28D9}{white}
\fm@curs{1batx-mi}{9B1C1C}{white}
\fm@curs{1batx-ccss}{EF6B6B}{white}
\fm@curs{2batx-mii}{1E6B47}{white}
\fm@curs{2batx-ccss}{2BB3A3}{white}

\newcommand{\fmcurs}[1]{%
  \@ifundefined{fm@color@#1}{\PackageError{fentmates}{Curs desconegut: #1}{}}{}%
  \edef\fm@tmp{\noexpand\definecolor{curs}{HTML}{\@nameuse{fm@color@#1}}%
               \noexpand\colorlet{curstext}{\@nameuse{fm@text@#1}}}\fm@tmp
  \colorlet{cursfosc}{curs!72!black}%
  \def\fm@logo{\fm@logodir/vertical-#1.pdf}%
  \def\fm@marca{\fm@logodir/marca-#1.pdf}}

% ---------- capçalera de la primera pàgina ----------
% #1 matèria (petit), #2 curs, #3 tipus de document i unitat, #4 títol
\newcommand{\capcalera}[4]{%
  \gdef\fm@cursnom{#2}%
  \par\noindent
  \begin{tcolorbox}[colback=curs,colframe=curs,boxrule=0pt,arc=2.2mm,width=3.5cm,height=2.15cm,
      valign=center,left=3mm,right=2mm,top=1mm,bottom=1mm,nobeforeafter,box align=center]
    \color{curstext}{\footnotesize #1}\par\vspace{1pt}{\fontsize{19}{21}\selectfont\bfseries #2\par}%
  \end{tcolorbox}\hspace{5mm}%
  \begin{minipage}[c]{\dimexpr\linewidth-3.5cm-5mm-2.9cm\relax}
    {\color{cursfosc}\small\bfseries\MakeUppercase{#3}}\par\vspace{2pt}
    {\color{navy}\fontsize{20}{23}\selectfont\bfseries #4\par}
  \end{minipage}\hfill
  \begin{minipage}[c]{2.6cm}\raggedleft\includegraphics[width=2.6cm]{\fm@logo}\end{minipage}\par
  \vspace{2pt}}

% ---------- barres de secció del color del curs ----------
\renewcommand{\seccio}[2]{\par\needspace{\fmneed\baselineskip}\bigskip
  \begin{tcolorbox}[colback=curs,colframe=curs,coltext=curstext,arc=2mm,boxrule=0pt,left=3mm,right=3mm,top=1.2mm,bottom=1.2mm]
  {\large\bfseries #1\quad #2}
  \end{tcolorbox}\vspace{-2pt}}

% ---------- estrelles de dificultat del color del curs (a Batxillerat no se'n posen) ----------
\@ifundefined{estrelles}{}{%
  \renewcommand{\estrelles}[1]{{\color{curs}\ifcase#1\or\faStar\faStar[regular]\faStar[regular]\or\faStar\faStar\faStar[regular]\or\faStar\faStar\faStar\fi}}}

% ---------- capçalera i peu de les pàgines ----------
% primera pàgina: logos institucionals al peu
\fancypagestyle{primera}{\fancyhf{}\renewcommand{\headrulewidth}{0pt}
  \fancyfoot[L]{\raisebox{-0.32\height}{\includegraphics[height=0.78cm]{logo-gva.png}}\hspace{5mm}%
    \raisebox{-0.32\height}{\includegraphics[height=0.78cm]{logo-vermellar.png}}}
  \fancyfoot[R]{\scriptsize\color{muted}\href{https://creativecommons.org/licenses/by-nc-sa/4.0/deed.ca}{\ccbyncsa}\ \ CC BY-NC-SA 4.0 · fentmates.org\qquad\small\thepage}}
% la resta: el símbol a l'esquerra i el curs a la dreta, amb el títol al centre com abans
\fancyhead[L]{\raisebox{-0.3\height}{\includegraphics[height=0.62cm]{\fm@marca}}}
\fancyhead[R]{\small\bfseries\color{cursfosc}\fm@cursnom}
\renewcommand{\headrule}{\color{curs!45}\hrule height \headrulewidth}
\makeatother

\fmcurs{1batx-mi}

\begin{document}
\thispagestyle{primera}

\capcalera{Matemáticas I}{1\textsuperscript{o} Bach.}{Ficha de ejercicios · Unidad 2}{\raggedright Álgebra. Ecuaciones e inecuaciones}
\vspace{10pt}
Nombre y apellidos: \rule{7.8cm}{0.4pt}\hfill Grupo: \rule{1.6cm}{0.4pt}\hfill Fecha: \rule{2.2cm}{0.4pt}
\vspace{8pt}

\begin{tcolorbox}[colback=soft,colframe=soft,arc=2mm,left=3mm,right=3mm,top=1.5mm,bottom=1.5mm]
\small Resuelve los ejercicios en el cuaderno mostrando todos los pasos. Da los resultados de manera exacta y, si hace falta, también aproximados. En las ecuaciones con denominadores, radicales o logaritmos, comprueba siempre las soluciones. Los ejercicios 1--62 son de la ficha del departamento; la sección 12 tiene otros nuevos para ampliar y aplicar lo que has aprendido.
\end{tcolorbox}

% =====================================================
\seccio{1}{Polinomios. Operaciones}

\exf{1} Simplifica las siguientes expresiones polinómicas.\begin{tasks}(1)
  \task $2(x-2)^2-3(3x+2)-2(3x-2)(3x+2)$
  \task $(3x+2)^2+2(2x-3)^2-(x-5)(2x-5)$
  \task $(2x^2-2x-13)(3x^2-2x)-3x$
  \task $(x^2+4)(x^2-4)(x-2)+2x(x^3-10)x$
  \task $\left(\frac23x-\frac35\right)\left(\frac32x^2+\frac25\right)+\frac{6}{25}$
\end{tasks}\recorda{$(a\pm b)^2=a^2\pm2ab+b^2$, \ $(a+b)(a-b)=a^2-b^2$. Desarrolla cada producto y reduce los términos semejantes.}

\ex{2} Demuestra la igualdad algebraica: \[(x+y+z)^2=x^2+y^2+z^2+2xy+2xz+2yz\]

\ex{1} Calcula $(4x^2-3x+2)^2$.

\ex{2} Calcula los valores de $a$ y $b$ para que se verifiquen las siguientes igualdades.\begin{tasks}(1)
  \task $(3x^2-5x+a)(2x^2+bx)=6x^4-19x^3+19x^2-6x$
  \task $2(x+3)^2-3(x+a)=2x^2+bx+3$
\end{tasks}

\ex{1} Efectúa las siguientes divisiones:\begin{tasks}(1)
  \task $(3x^3+2x^2+x-5):(3x^2+2)$
  \task $(3x^4-2x^2-x+4):(x^2+x+2)$
  \task $(x^3-3x^2-x+6):(2x^3-2x+3)$
\end{tasks}

\ex{1} Efectúa las siguientes divisiones aplicando la regla de Ruffini:\begin{tasks}(2)
  \task $(3x^3+2x^2+x-5):(x+2)$
  \task $(3x^4-2x^2-x+4):(x+2)$
  \task $(x^3-3x^2-x+6):(2x+3)$
  \task $(4x^4-3x^2+5):(-x+5)$
  \task $(2x^5-12x^3-6x^2-117):(x-3)$
\end{tasks}\recorda{Si el divisor es $ax+b$, haz Ruffini con $-\frac ba$ y divide el cociente entre $a$; el resto no cambia.}

\ex{2} Utiliza Ruffini para calcular el valor de $a$ y $b$ que hace que las siguientes divisiones sean exactas:\begin{tasks}(1)
  \task $(2x^3+5x^2-2x+a):(x+3)$
  \task $(2x^4+3x^3-2x^2-ax+3):(2x+3)$
  \task $(-6x^4+11x^3+ax^2+5x+4):(-3x+4)$
  \task $(4x^4+6x^3-13x^2+x+b):(-x+5)$
\end{tasks}

\ex{1} Comprueba si $M(x)=2x^3-5x^2+4x-4$ es divisible por $x-2$ y, si lo es, halla un polinomio $N(x)$ que permita escribir $M(x)=(x-2)\cdot N(x)$.


% =====================================================
\seccio{2}{Factorización de un polinomio. Teorema del resto}

\exf{1} Comprueba si son raíces del polinomio $P(x)=x^4+3x^3-2x^2+6x-8$ los números siguientes:\begin{tasks}(5)
  \task $x=1$
  \task $x=2$
  \task $x=-1$
  \task $x=-4$
  \task $x=-3$
\end{tasks}

\ex{1} Calcula el valor numérico del polinomio $P(x)=x^2-2x-3$ para $x=-1$ y $x=3$. ¿Podrías encontrar otro valor de $x$ para el cual el valor del polinomio sea cero? Razona la respuesta.

\ex{1} Determina las raíces de los polinomios siguientes:\begin{tasks}(2)
  \task $(x-3)(x+5)(x-2)$
  \task $x(x-2)^3(2x+1)$
  \task $(2x-1)(3x+2)(x+1)^3$
  \task $x^3-3x^2-6x+8$
  \task $x^3+8x^2+17x+10$
  \task $2x^5-11x^4+21x^3-16x^2+4x$
\end{tasks}

\ex{2} Determina, si existen, las raíces enteras de cada uno de los siguientes polinomios y factorízalos.\begin{tasks}(2)
  \task $P(x)=x^3-2x^2-5x+6$
  \task $P(x)=x^3+x^2-5x+3$
  \task $P(x)=2x^4-x^3-5x^2+x+3$
  \task $P(x)=4x^3-4x^2-11x+6$
  \task $P(x)=x^4-4x^2$
  \task $P(x)=6x^5+11x^4+3x^3-3x^2-x$
  \task $P(x)=x^6-4x^4-x^2+4$
  \task $P(x)=\frac{x^4}{4}+\frac34x^2+\frac{9}{16}$ \ (comprueba que no tiene raíces reales)
\end{tasks}\recorda{Las raíces enteras son divisores del término independiente. Pruébalas con Ruffini hasta que quede un polinomio de grado 2.}

\ex{1} Calcula el valor de $k$ para que el polinomio:\begin{tasks}(1)
  \task $P(x)=x^3+x^2-2x+k$ sea divisible por $x-2$
  \task $P(x)=x^3-2x^2+kx+4$ sea divisible por $x+2$
  \task $P(x)=(k-3)x^3-15x+k$ sea divisible por $x+4$
  \task $P(x)=x^3-(k-3)x^2+(k-5)x-3$ sea divisible por $x-3$
\end{tasks}

\ex{2} Calcula el valor de $k$ para que la división de $x^3-(3k+2)x^2-3$ entre $(x+3)$:\begin{tasks}(2)
  \task Sea exacta.
  \task Tenga resto 15.
\end{tasks}

\ex{2} De un polinomio de segundo grado, $P(x)$, sabemos que $P(1)=-6$, $P(0)=-3$ y una de sus raíces es 3. Determínalo.

\ex{1} Halla el valor de $m$ para que el polinomio $mx^3-6x^2-4x+8$ tenga 2 como raíz.

\ex{2} Determina $a$ y $b$ de manera que el polinomio $x^3+ax^2+bx-6$ sea divisible por $x-2$ y por $x+3$.

\ex{1} Halla el valor de $m$ para que el polinomio $2x^4+9x^3+2x^2-6x+3m$ tenga resto 12 al dividirlo por $x+2$, sin hacer la división.


% =====================================================
\seccio{3}{Fracciones algebraicas}

\exf{2} Simplifica las siguientes fracciones algebraicas.\begin{tasks}(2)
  \task $\dfrac{2x^2-8}{x^3+x^2-2x}$
  \task $\dfrac{x^3-x}{x^3+3x^2+2x}$
  \task $\dfrac{x^4-2x^3-3x^2}{x^4-9x^2}$
  \task $\dfrac{x^3-3x^2+3x-1}{x^3-2x^2+x}$
  \task $\dfrac{2x^4-5x^3+x^2+5x-3}{2x^3-x^2-7x+6}$
  \task $\dfrac{2x^4+x^3-11x^2+11x-3}{2x^3+3x^2-8x+3}$
  \task $\dfrac{x^3-2x^2-9x+18}{x^3-7x^2+16x-12}$
  \task $\dfrac{x^4-5x^3+9x^2-7x+2}{x^4-6x^3+13x^2-12x+4}$
\end{tasks}\recorda{Factoriza el numerador y el denominador y simplifica solo los \emph{factores} comunes.}

\ex{2} Efectúa las siguientes sumas y restas reduciendo previamente a común denominador:\begin{tasks}(2)
  \task $\dfrac{1}{x+1}+\dfrac{3}{x^2-1}$
  \task $\dfrac{1}{x^2+x}+\dfrac{3}{x+1}$
  \task $\dfrac{1}{2x}-\dfrac{3}{4x^2-4x}$
  \task $\dfrac{1}{x^2+x}+\dfrac{3}{x^2-x}$
  \task $\dfrac{2x}{x-2}+\dfrac{3x-1}{x+2}-\dfrac{1}{x^2-4}$
  \task $\dfrac{1}{(x-1)^2}+\dfrac{2}{x-1}+\dfrac{1}{x^2-1}$
\end{tasks}

\ex{2} Opera y simplifica las fracciones algebraicas:\begin{tasks}(2)
  \task $\dfrac{3x+3}{12x-12}:\dfrac{(x+1)^2}{x^2-1}$
  \task $\dfrac{x^2+2x-3}{(x-2)^3}\cdot\dfrac{(x-2)^2}{x^2-1}$
  \task $\dfrac{6x+4}{x-2}\cdot\dfrac{4x+1}{2x+4}$
  \task $\dfrac{-6x^2}{x^2+2x+1}\cdot\dfrac{x+1}{3x^3}\cdot\dfrac{x+1}{3x}$
  \task $\left(\dfrac3x-\dfrac{2x}{x+1}\right)\cdot\dfrac{x^2+x}{x-1}$
  \task $\dfrac{(x-1)^2}{2}\cdot\dfrac{1}{x^2-1}-\dfrac{3x}{(x+1)^2}$
  \task $\left(\dfrac{2x-1}{x+1}-\dfrac{3x}{x-1}\right)\cdot\dfrac{x^3-x}{-x^2-6x+1}$
\end{tasks}

\ex{3} Calcula y simplifica el resultado de las operaciones siguientes:\begin{tasks}(2)
  \task $\dfrac{x+1}{x}+\dfrac{2x+1}{x+2}-\dfrac{x}{x^2+1}$
  \task $\dfrac{1+\frac{2x+1}{x-1}}{\frac{1}{x-1}+\frac{x}{x+1}}$
  \task $\dfrac{ba-a^2}{a^3b-ab^3}:\dfrac{b}{a^2+2ab+b^2}$
  \task $\dfrac{a^2}{ab}+\dfrac{ab^2}{b^4}-a$
  \task $\dfrac{6}{2+x}-\dfrac{4}{2-x}+\dfrac{16}{x^2-4}$
  \task $\dfrac{2x^2-7x+6}{x^2+2x+1}\cdot\dfrac{x^2+6x+5}{2x-3}$
  \task $(x^2-y^2):\left(\dfrac1x+\dfrac1y\right)$
  \task $\dfrac{a+x}{x^2-a^2}\cdot\dfrac{x-a}{x+a}$
  \task $1+\cfrac{1}{1+\cfrac{1}{1+\frac1x}}$
  \task $1+\dfrac{\frac{2x+1}{x-1}}{\frac{1}{x-1}+\frac{x}{x+1}}$
\end{tasks}


% =====================================================
\seccio{4}{Ecuaciones polinómicas}

\exf{2} Resuelve las siguientes ecuaciones polinómicas:\begin{tasks}(2)
  \task $\dfrac{2x-3}{4}+\dfrac x2-\dfrac{3x-1}{5}=2x-1$
  \task $2(3x-2)+x(x-1)=-4$
  \task $(x+1)^3-(x-1)^3=7$
  \task $\dfrac{x^2+1}{2}-\dfrac{2x-3}{4}+\dfrac{x^2}{6}=\dfrac{59}{12}$
  \task $6x^4+13x^3-8x^2-17x+6=0$
  \task $2x^4-x^3-3x-18=0$
  \task $6x^4+5x^3=2x^2+x$
  \task $x^3-7x+6=0$
  \task $x^3-3x^2-9x+27=0$
  \task $x^4-x^3-5x^2-x-6=0$
  \task $x^4-1=0$
  \task $x^4+x^3-9x^2-9x=0$
\end{tasks}\recorda{Pásalo todo a un miembro y factoriza (factor común, Ruffini…): un producto vale cero si algún factor vale cero.}

\ex{2} Resuelve las siguientes ecuaciones bicuadradas (o que se resuelven con un cambio de variable):\begin{tasks}(2)
  \task $x^4-17x^2+16=0$
  \task $x^6-9x^3+8=0$
  \task $x^4-8x^2-9=0$
  \task $4x^6+8x^3-12=0$
  \task $2(x+1)^4-8x^3-8(x+3)+8=0$
  \task $\dfrac{9}{x^2}=10-x^2$
\end{tasks}

\ex{1} Resuelve las siguientes ecuaciones factorizadas:\begin{tasks}(2)
  \task $(x^2-3)(x^2-4)=0$
  \task $(3x+2)(x^2-4x)(x^2-9)=0$
  \task $(x^2-4x+4)(x^2+6x+9)=0$
  \task $(x-2)^2(x+1)^3(2x-7)=0$
\end{tasks}


% =====================================================
\seccio{5}{Ecuaciones racionales}

\exf{2} Resuelve las ecuaciones racionales siguientes:\begin{tasks}(2)
  \task $\dfrac1x+\dfrac1{x^2}+\dfrac1{x^3}=\dfrac78$
  \task $\dfrac{2x}{3}+\dfrac{2x+3}{x-1}=\dfrac{11}{3x-3}$
  \task $\dfrac{2x}{x-2}+\dfrac{3x}{x+2}=\dfrac{6x^2}{x^2-4}$
  \task $\dfrac{3}{x^2-6x+8}+\dfrac{x-5}{x^2-5x+4}=\dfrac{1}{x-4}$
  \task $\dfrac{x}{\sqrt2}+\dfrac{\sqrt2}{x}=\sqrt2\,x$
  \task $\dfrac{x}{x-2}+\dfrac{2x}{2-x}=-x$
  \task $\dfrac{1+\frac{x+1}{x-1}}{2-\frac{x-1}{x+1}}=2$
  \task $\dfrac{1}{1-\frac{1}{x+1}}=\dfrac{x+1}{x}$
\end{tasks}\recorda{Multiplica por el m.c.m. de los denominadores y, al final, descarta las soluciones que anulan algún denominador.}


% =====================================================
\seccio{6}{Ecuaciones con radicales}

\exf{2} Resuelve las siguientes ecuaciones con radicales:\begin{tasks}(2)
  \task $-\sqrt{2x-3}+1=x$
  \task $x+\sqrt{2x^2+2x-3}=-1$
  \task $\sqrt{\dfrac{7x+1}{4}}=\dfrac{5x-7}{6}$
  \task $\dfrac{\sqrt{2x-1}}{4}=\dfrac{3}{\sqrt{2x-1}}$
  \task $\sqrt{2-5x}+x\sqrt3=0$
  \task $\sqrt{2+\sqrt{x-4}}=\sqrt{12-x}$
  \task $\sqrt x+\sqrt{x-4}=2$
  \task $\sqrt{x+5}+\sqrt{2x+8}=7$
  \task $\sqrt{x+4}+\sqrt{x-1}=5$
  \task $\sqrt{x-7}+\sqrt{2x}=\sqrt{x+1}$
  \task $\sqrt{1-4x}-2\sqrt{3+x}=\sqrt{3+x}$
\end{tasks}\recorda{Aísla un radical, eleva al cuadrado y repite si hace falta. Comprueba \emph{siempre} las soluciones en la ecuación inicial.}

\ex{2} Resuelve las siguientes ecuaciones con radicales no cuadráticos:\begin{tasks}(2)
  \task $\sqrt[3]{x^2-1}+1=x$
  \task $\dfrac{\sqrt{3x+7}}{\sqrt[3]{x+1}}=\sqrt[3]{x^2+2x+1}$
  \task $\sqrt[3]{4x-1}=x-4$
  \task $\sqrt[3]{4-2x}=\sqrt[6]{8x^2-16x}$
  \task $\sqrt{2x+2}-\sqrt[4]{6x+10}=0$
\end{tasks}\recorda{Eleva los dos miembros al índice de la raíz (o al m.c.m. de los índices).}

\ex{2} Calcula el valor de un número tal que, si se le suma una unidad y después se extrae la raíz cuadrada, se obtiene el doble que al restarle 11 unidades y extraer la raíz cuadrada.


% =====================================================
\seccio{7}{Ecuaciones logarítmicas}

\exf{2} Resuelve las siguientes ecuaciones logarítmicas:\begin{tasks}(2)
  \task $\log2+\log(x+3)=\log(x+5)$
  \task $\log(15-2x)=2\log x$
  \task $\log(2x+3)-\log(x-1)=2\log2+2\log3$
  \task $\frac12\log(x^2-3x+2)=\log x$
  \task $\log\left(\dfrac{2x-2}{x}\right)=2\log(x-1)-\log x$
  \task $2\log x=\log\dfrac x2-1$
  \task $\log(4-5x)+\log(2x)=\log(2x-x^2)+\log3$
  \task $\log_2(5x^2+15x+10)-\log_2(x+2)=2$
  \task $\dfrac{\log(4-x)}{\log(x+2)}=2$
  \task $\log_{x+3}9=\log_53$
  \task $\log_332+\log_{\frac13}(6-x)=\log_{\sqrt3}x$
\end{tasks}\recorda{Agrupa los logaritmos con las propiedades hasta llegar a $\log A=\log B$ (y entonces $A=B$). Comprueba que todos los argumentos son positivos.}

\ex{1} Halla el número que cumple que, al añadir a su logaritmo decimal el logaritmo decimal de 2, el resultado es 1.

\ex{1} Calcula el valor de un número sabiendo que el doble de su logaritmo decimal es igual a la suma de los logaritmos decimales de 4 y de 9.


% =====================================================
\seccio{8}{Ecuaciones exponenciales}

\exf{2} Resuelve las siguientes ecuaciones exponenciales:\begin{tasks}(2)
  \task $4^{x^2+1}=2^{5x+5}$
  \task $4^{(x-2)^2}=262144$
  \task $3^{x+4}=5$
  \task $\dfrac{1}{2^x}=16^{\frac{x(x-1)}{2}}$
  \task $13^{x^2+2x}-\dfrac1{13}=0$
  \task $2^{2x^2-4x-2}=16$
  \task $2\cdot6^{4x-3}=24$
  \task $3\cdot2^x=2\cdot3^x$
  \task $\sqrt[x]{27}=3^{x+2}$
  \task $2^{x+4}-8^x=0$
  \task $7^{x^2-5x+6}=1$
  \task $49^{\frac{x-5}{3}}=\dfrac{1}{\sqrt{7^{x-1}}}$
\end{tasks}\recorda{Si puedes, escribe los dos miembros con la misma base e iguala los exponentes; si no, toma logaritmos.}

\ex{2} Resuelve las siguientes ecuaciones exponenciales:\begin{tasks}(2)
  \task $2^{x+1}+2^x+2^{x-1}=28$
  \task $5^{x+3}-5^{x-1}-3120=0$
  \task $4^x+4^{x-1}+4^{x-2}=336$
  \task $3^{x+2}+3^{x+1}+3^x+3^{x-1}=120$
  \task $2^{x-1}+2^{x-2}+2^{x-3}+2^{x-4}=960$
  \task $6^x-9\cdot6^{-x}+8=0$
\end{tasks}\recorda{Descompón las potencias ($a^{x+1}=a\cdot a^x$, $a^{x-1}=\frac{a^x}{a}$) y saca factor común $a^x$.}

\ex{3} Resuelve las siguientes ecuaciones exponenciales:\begin{tasks}(2)
  \task $3^{2x}-3^{x-1}=3^{x+1}-1$
  \task $2\cdot10^{2x+4}+3\cdot10^{x+2}-5=0$
  \task $9^x+5\cdot3^x-24=0$
  \task $9^x+3^{2x-1}+3^{x-1}=111$
  \task $3^{x+2}+9^{x-1}=90$
  \task $2^{2x-4}-5\cdot2^{x-3}+1=0$
  \task $9^{x+2}+3^{x+3}-810=0$
  \task $4^{x+1}+2^{x+3}=320$
\end{tasks}\recorda{Si hay $a^{2x}$ y $a^x$, haz el cambio $t=a^x$; las soluciones $t\le0$ no valen.}


% =====================================================
\seccio{9}{Problemas de ecuaciones}

\exf{2} Un comerciante tiene un tonel de vino de la clase A de precio 5\,€/L. Añade 50 litros de vino de la clase B de 2\,€/L y 50 litros de agua y la mezcla resulta a 4\,€/L. ¿Cuántos litros de la clase A tenía el tonel?

\ex{1} Tenemos 30 litros de vino a 5\,€/L y queremos añadirle vino de 10\,€/L para obtener una mezcla que salga a 7\,€/L. ¿Cuántos litros tenemos que añadir?

\ex{2} En una tienda se vende té blanco a 18\,€/kg y té verde a 14\,€/kg. También se vende una mezcla de los dos a 16,4\,€/kg. ¿Cuál es la composición de la mezcla?

\ex{1} En un colegio hay tres cursos. El 60\,\% de los alumnos estudia primero, el 30\,\% estudia segundo y el resto, que son 15 alumnos, estudia tercero. ¿Cuántos alumnos hay en el colegio?

\ex{2} El 80\,\% de las entradas de un partido se reservaron para los socios del club; de las restantes, el 25\,\% se enviaron a los socios del club contrario y las demás, 18\,000, se pusieron a la venta. ¿Cuántas entradas había?

\ex{1} Una fábrica exporta el 30\,\% de la producción a Asia, el 40\,\% a Europa, el 15\,\% a América y el resto, que supone una facturación de 3 millones de euros, a África. ¿Cuál es la facturación total?

\ex{2} Un examen consta de 20 preguntas. Cada acierto suma 5 puntos y cada error resta 2 puntos. Una alumna sacó 58 puntos. ¿Cuántas preguntas acertó y cuántas falló, si no se puede dejar ninguna sin contestar?

\ex{3} Un tendero invierte 125\,€ en la compra de una partida de manzanas. Desecha 20 kilos por defectuosas y vende el resto por 147\,€, ganando 0,40\,€ por kilo sobre el precio de compra. ¿Cuántos kilos compró?

\ex{1} Calcula la base y la altura de un triángulo rectángulo si son números enteros consecutivos y la hipotenusa mide 5 cm.

\ex{2} Calcula la longitud de los tres lados de un triángulo rectángulo si son números pares consecutivos.

\ex{2} Calcula las dimensiones de una finca rectangular si el perímetro mide 140 m y la diagonal 50 m.

\ex{1} Una balsa tarda 5 horas en llenarse con la toma de agua habitual y 20 horas con una manguera. ¿Cuánto tiempo hace falta si usamos los dos métodos a la vez?

\ex{1} Un hombre tarda 2 horas en barrer una casa y otro tarda 3. ¿Cuánto tardarían trabajando juntos?

\ex{2} Dos personas trabajando juntas tardan 10 horas en hacer un trabajo, pero una trabaja el doble de rápido que la otra. ¿Cuánto tardaría cada una por separado?

\ex{2} Una piscina se puede llenar con 5 fuentes que, por separado, tardan 10, 20, 30, 40 y 50 horas respectivamente. ¿Cuánto tardarán las cinco juntas?

\ex{3} Un grupo de personas contrata un autobús por 600\,€ a pagar a partes iguales; al final 10 personas no pueden ir al viaje y cada persona que va tiene que pagar 5\,€ más. ¿Cuántas personas iban inicialmente?

\ex{2} El número de enfermos de gripe en una ciudad a lo largo del pasado enero viene dado por $y(t)=100+200\cdot e^{0,2t}$, donde $t$ son los días transcurridos desde el 1 de enero.\begin{tasks}(1)
  \task ¿Cuántos enfermos había el 1 de enero?
  \task ¿Qué día el número de enfermos era de 900?
\end{tasks}

\ex{2} La cantidad que queda de una sustancia radiactiva después de $t$ días es $q(t)=M\cdot e^{-0,1t}$.\begin{tasks}(1)
  \task ¿Después de cuánto tiempo la masa se habrá reducido a la mitad?
  \task Si la masa inicial es de 27 mg, ¿cuánta sustancia quedará después de 10 días?
\end{tasks}

\ex{2} Si inicialmente hay $n$ bacterias en un cultivo y el número se duplica cada hora, después de $t$ horas hay $N=n\cdot2^t$. Si inicialmente hay 4500, ¿cuánto tiempo tardará el cultivo en llegar a 50\,000 bacterias?

\ex{3} La biomasa $C$ (en kg) por unidad de superficie de una región sigue la función $C(x)=\dfrac{4}{1+K\cdot a^{x/10}}$, con $x\ge0$ el tiempo en años.\begin{tasks}(1)
  \task Calcula $K$ y $a$ para que a los 10 años haya 1 kg de biomasa y a los 20 años, 1,6 kg.
  \task ¿Cuántos años tienen que pasar para que la biomasa sea de 3 kg?
\end{tasks}


% =====================================================
\seccio{10}{Inecuaciones con una incógnita}

\exf{1} Comprueba si los números reales indicados pertenecen a la solución de las inecuaciones correspondientes:\begin{tasks}(1)
  \task $x=-2$ de la inecuación $x^3+x^2+x\le6$
  \task $x=\dfrac1e$ de la inecuación $x+\ln x\ge0$
  \task $x=-0,5$ de la inecuación $2^x+x-2\lt3^x$
\end{tasks}

\ex{1} Resuelve las siguientes inecuaciones:\begin{tasks}(2)
  \task $\dfrac{x-3}{2}-\dfrac{x-2}{8}\le\dfrac x2$
  \task $2x-3-\dfrac x2\gt x+\dfrac{3x+1}{6}$
  \task $x+2(x+1)+3(x+2)\lt\dfrac{x+38}{2}$
  \task $\dfrac{x^2-1}{2}-\dfrac{3x-2}{3}\le\dfrac{(x+1)^2}{6}$
\end{tasks}

\ex{2} Resuelve las siguientes inecuaciones:\begin{tasks}(2)
  \task $-2x^3+3x\gt0$
  \task $3x^2+x-24\ge0$
  \task $x^3-4x\gt0$
  \task $x^3+x^2-x-1\le0$
  \task $3x(2x-1)+x^2\ge5x-1$
  \task $-x^3-2x^2+x+2\le0$
  \task $x^4-1\gt0$
  \task $(x-1)(x+4)\lt-6$
\end{tasks}\recorda{Pásalo todo a un miembro, factoriza y haz la tabla de signos. Los extremos se incluyen si la desigualdad es $\le$ o $\ge$.}

\ex{2} Resuelve las inecuaciones racionales siguientes:\begin{tasks}(2)
  \task $\dfrac{5x-2}{2x+1}\ge0$
  \task $\dfrac{x^2-x-2}{x^2+x-2}\ge0$
  \task $\dfrac{4x-5}{4x^2-x-5}\lt0$
  \task $\dfrac{x^3}{x-9}\ge0$
  \task $\dfrac{x^2+4x-5}{x^3-x}\ge0$
  \task $\dfrac{2x+3}{x^3+1}\ge0$
\end{tasks}\recorda{Las raíces del denominador no se incluyen nunca en la solución.}


% =====================================================
\seccio{11}{Problemas de inecuaciones}

\exf{1} Una compañía telefónica ofrece dos tarifas: la A tiene una cuota fija de 10\,€ al mes más 0,05\,€ por minuto, y la B no tiene cuota fija pero cobra 0,12\,€ por minuto. ¿A partir de cuántos minutos mensuales es más económica la tarifa A?

\ex{1} La nota final de una asignatura se calcula con el 30\,\% de la nota del primer examen y el 70\,\% de la del segundo. Si una alumna ha sacado un 4 en el primero, ¿qué nota tiene que sacar en el segundo para aprobar (nota final $\ge5$)? Las notas van de 0 a 10.

\ex{2} Se lanza una pelota hacia arriba y su altura, en metros, es $h(t)=20t-5t^2$, con $t$ en segundos $(0\le t\le4)$. ¿Durante cuánto tiempo está la pelota por encima de los 15 metros?


% =====================================================
\seccio{12}{Ampliación y problemas}

\exf{1} Calcula, sin hacer la división, el resto de $(x^{100}-2x^{51}+1):(x+1)$.

\ex{2} Halla el polinomio de grado 3 que tiene las raíces $1$, $-2$ y $3$ y que cumple $P(0)=12$.

\ex{2} Simplifica:\begin{tasks}(2)
  \task $\dfrac{x^3-27}{x^2-9}$
  \task $\dfrac{1-\frac1x}{1-\frac1{x^2}}$
\end{tasks}

\ex{2} Resuelve $\dfrac{x}{x+1}+\dfrac{x+1}{x}=\dfrac52$.

\ex{2} Resuelve $\sqrt{x+3}-\sqrt{x-2}=1$.

\ex{2} Resuelve $2^{2x+1}-5\cdot2^x+2=0$.

\ex{2} Resuelve $\log_2(x+1)+\log_2(x-1)=3$.

\ex{2} Resuelve $\dfrac{x^2-4}{x^2+1}\le0$. ¿Por qué no hace falta estudiar el signo del denominador?

\ex{2} Ingresamos 5000\,€ en una cuenta al 4\,\% de interés compuesto anual. ¿Cuántos años completos tienen que pasar para que haya al menos 8000\,€?

\vfill
\begin{center}\color{muted}\small\faCheckCircle[regular]\ \ ¡Ánimo! Comprueba tus resultados con las herramientas interactivas de los apuntes de la unidad.\end{center}

\end{document}
